\paragraph{Persistent assignment artifacts.}
Assignment preparation is also demand-independent. Its immutable reusable state
consists of the time-expanded graph arrays and labels, destination/OD grouping
arrays, and generalized-cost components. It does not contain demand magnitudes,
observations, priors, optimizer settings, or compiled executables. The package
can store this state as compressed NPZ plus canonical JSON through optional
\texttt{off}, \texttt{auto}, \texttt{refresh}, and \texttt{readonly} policies;
ordinary uncached calls remain the default.

The assignment-cache provenance includes the complete graph-relevant scenario
contents, OD keys, assignment and cost configuration, numeric type, package and
schema versions, sentinel conventions, and indexing rules. Loaded archives are
checked for exact provenance, recognized fields, counts, shapes, dtypes, CSR and
padded-mask consistency, and integer/Boolean conventions before JAX arrays are
reconstructed. Cache writes use a unique temporary file, synchronization to
disk, and atomic replacement. Thus concurrent writers are safe, corrupt or
incompatible entries rebuild under \texttt{auto}, and \texttt{readonly} never
writes. The measurement-operator cache retains its own compact-input content
hash, so a valid assignment archive cannot validate an incompatible operator.

Internal TPG profiling initially assigned 3.142 s of a 3.313 s preparation to
destination grouping. The cause was repeated Python scalar extraction from JAX
node arrays, which introduced thousands of small device synchronizations. One
NumPy view is now reused for immutable node indexing and validation, and all
destination link masks are constructed in one broadcasted operation. Uncached
assignment construction consequently fell to 0.097 s: graph construction took
about 0.049 s, OD indexing and grouping 0.022 s, and cost construction plus
synchronization 0.026 s in the compilation-cache-enabled run.

The TPG artifact contains 36,478 nodes, 11,259 links, 22,022 OD records, and 113
destination groups. Its explicit arrays occupy 6,909,418 bytes and compress to
576,227 bytes. A representative strict cache hit spent 0.00044 s opening the
archive, 0.00460 s decompressing arrays, 0.01473 s in combined validation and
decompression, 0.000006 s reconstructing host objects, 0.00500 s transferring
and synchronizing JAX arrays, and 0.08787 s generating the canonical provenance.
The complete assignment-cache call took 0.1084 s.

\begin{table}[htbp]
\centering
\small
\begin{tabular}{lrrrrr}
\hline
Policy/case & Assignment & Preparation & Warm V\&G & Optimizer & Process \\
\hline
Off & 0.0971 & 1.1756 & 0.000992 & 0.2908 & 2.1939 \\
Populate & 0.3412 & 1.3532 & 0.000993 & 0.2975 & 2.4106 \\
Valid hit & 0.1084 & 1.0775 & 0.000937 & 0.2721 & 2.0211 \\
Corrupt/rebuild & 0.3201 & 1.3339 & 0.000965 & 0.3073 & 2.4081 \\
Mismatch/rebuild & 0.3321 & 1.3463 & 0.001048 & 0.3166 & 2.3935 \\
Read-only hit & 0.1128 & 1.1242 & 0.000972 & 0.2922 & 2.1872 \\
\hline
\end{tabular}
\caption{Fresh-process TPG assignment-cache measurements in seconds with the
measurement-operator and JAX compilation caches enabled.}
\label{tab:tpg-assignment-cache}
\end{table}

After eliminating scalar synchronization, uncached reconstruction is already
slightly faster than strict cache validation in the assignment stage itself;
canonical fingerprinting accounts for most of the hit time. The persistent
artifact is therefore useful primarily for strict reproducibility and validation,
while the focused grouping simplification provides the main speedup. Relative to
the original profile, complete one-iteration time falls from 5.389 s to roughly
2.0--2.2 s. Cache-hit runs at 1, 5, and 20 iterations reproduce objectives
55,243.4375, 49,635.4922, and 49,140.4258 exactly.

